\documentclass[11pt,aps,preprint,nofootinbib,superscriptaddress,tightenlines]{revtex4-2}
\usepackage[T1]{fontenc}
\usepackage{amsmath,amssymb,mathtools}
\let\lambdabar\relax   % revtex defines it; newpxmath does too
\usepackage{newpxtext,newpxmath}
\usepackage{bm}
\usepackage[final]{microtype}
\usepackage{xcolor}
\usepackage{booktabs}
\usepackage{graphicx}
\usepackage{dsfont}
\usepackage[most]{tcolorbox}
\usepackage{algpseudocode}
\newcommand{\one}{\mathds{1}}
\newcommand{\avg}[1]{\langle #1\rangle}
\newcommand{\klab}{\mathcal B}
\definecolor{darkblue}{RGB}{20,60,140}
\definecolor{accent}{RGB}{30,80,150}
\definecolor{tint}{RGB}{243,246,251}
\newtcolorbox{algobox}[1]{enhanced, colback=tint, colframe=accent, boxrule=0.6pt, arc=2pt,
  left=6pt, right=6pt, top=4pt, bottom=4pt, fonttitle=\bfseries\small, coltitle=white,
  colbacktitle=accent, title={#1}, before skip=10pt, after skip=10pt}
\newtcolorbox{keyeq}{enhanced, colback=tint, frame hidden, borderline west={2pt}{0pt}{accent},
  sharp corners, left=4pt, right=4pt, top=0pt, bottom=0pt, before skip=6pt, after skip=6pt}
\usepackage[colorlinks=true,linkcolor=darkblue,citecolor=darkblue,urlcolor=darkblue]{hyperref}

\def\st{\begin{equation}}
\def\stp{\end{equation}}
\def\Eq#1{eq.~(\ref{#1})}
\def\Sect#1{Sect.~\ref{#1}}
\def\Fig#1{Fig.~\ref{#1}}
\def\App#1{Appendix~\ref{#1}}
\def\tr{\operatorname{tr}}
\let\oldparagraph\paragraph
\renewcommand{\paragraph}{\par\bigskip\oldparagraph}


\begin{document}
\sloppy

\title{The split renormalization (Step 5) for any constant speed of sound:\\ development plan}
\author{D.~Teaney}
\affiliation{Department of Physics and Astronomy, Stony Brook University,
Stony Brook, New York 11794, USA}
\date{\today}

\begin{abstract}
A plan for implementing the renormalized particle Hamiltonian of Step 5 (\nolinkurl{partition.tex}, Appendix ``The label $s_{\rm eff}$''): the part of the
one-loop counterterm that depends on the label alone stays in the measure, as in \texttt{SPH\_RENORM=1}, and the rest, the non-conformal and the sound terms, goes
into the Hamiltonian.  It reproduces the target to $O(T/s_a)$ per particle for any constant $c_s^2$, and in the single-particle quadrature its residuals are $5$--$60$
times smaller than those of the Legendre Hamiltonian (\nolinkurl{legendre_plan.tex}).  Written on 2026-11-23, before any change to \texttt{src/}.
\end{abstract}

\maketitle
\tableofcontents

%=============================================================================
\section{The goal}
\label{sec:goal}

At one loop the canonical ensemble of the tree-level Hamiltonian has the target thermodynamics plus terms of order $T$ per particle: the label gas and each
particle's fluctuations, $\delta\ln Z_1$ per particle.  A renormalized model removes them.  Two such models exist:
\begin{itemize}
\item \texttt{SPH\_RENORM=1} removes the part that depends on the label alone, by the measure.  It is exact for the conformal fluid at rest, but it leaves the
  sound term in motion (the violation $R$ of the density-frame relation) and, for $c_s^2\neq1/2$, a pressure error $\mathcal N/\sigma$.
\item \texttt{SPH\_RENORM=legendre} puts the whole counterterm into the Hamiltonian by a Legendre transform.  It removes everything at one loop, but its
  residual is ten times larger, every map needs a two-dimensional Newton solve, and $H$ is two-valued at very negative labels.
\end{itemize}
Step 5 keeps the measure of \texttt{SPH\_RENORM=1} for the label part and adds the rest to the Hamiltonian as an energy.  The aim of this plan: implement it as
\texttt{SPH\_RENORM=split}, test it, time it, and repeat the measurements of the campaign \texttt{eos\_legendre} with it.

%=============================================================================
\section{For the NERSC session: working beside the running campaign}
\label{sec:nersc}

The campaign \texttt{eos\_legendre} finished on 2026-11-23 (53d29e4, the branch point of \texttt{split}; \nolinkurl{nersc/LOG.md}).  Its worktrees
(\texttt{\textasciitilde/stochastic-sph-legV} and others, each at a fixed commit) and its output stay as they are, for reruns and the comparison.  The rules:
\begin{enumerate}
\item \emph{Branch.}  \texttt{split} is made from \texttt{legendre} (this note is its first commit).  Work and push only on \texttt{split}; do not commit to
  \texttt{legendre}.  When \texttt{legendre} gets a fix to the shared code (\nolinkurl{src/renorm.jl}, \nolinkurl{src/pointwise.jl}, \nolinkurl{src/stochastic.jl}),
  merge \texttt{legendre} into \texttt{split}, never the reverse.
\item \emph{Checkout.}  A new worktree, \texttt{git worktree add \textasciitilde/stochastic-sph-split split}; precompile there with \texttt{julia --project=nersc}.
  Do not change the branch or the files of any existing worktree or of \texttt{\textasciitilde/stochastic-sph}.
\item \emph{Files.}  New files only, or files the \texttt{legendre} session does not edit: \nolinkurl{src/renorm_split.jl}, \nolinkurl{test/testsplit.jl},
  \nolinkurl{nersc/campaign_eos_split.txt}, \nolinkurl{examples/eos_split.py} (a copy of \nolinkurl{examples/eos_legendre.py}), figures \nolinkurl{docs/figs/split_*},
  and the section ``Results'' of this note.  Never edit \nolinkurl{legendre_plan.tex} or \nolinkurl{docs/figs/legendre_*}.  In \nolinkurl{nersc/LOG.md}, add a
  separate section ``\texttt{eos\_split}'' and append only to it, so that the later merge is a concatenation.
\item \emph{Output.}  \texttt{\$SCRATCH/sph/eos\_split/}, its own campaign directory.
\item \emph{Fairness of the comparison.}  The \texttt{legendre} code path must not change: with \texttt{SPH\_RENORM=legendre}, and with the switch off, the regression
  on \texttt{split} is bit-identical to \texttt{legendre} at the branch point.
\item \emph{Queue and budget.}  Implementation and tests in \texttt{-q debug} at any time.  Submit the campaign (about $1.1$ node-hours, \Sect{sec:work}F) when no
  \texttt{eos\_legendre} job is left in the queue (\texttt{squeue --me}); record the estimate in the log first.
\item \emph{The side-by-side.}  When both are done: \nolinkurl{examples/eos_compare.py}, reading \nolinkurl{docs/figs/legendre_table.csv} and
  \nolinkurl{docs/figs/split_table.csv}, the same panels with the tree level, \texttt{SPH\_RENORM=1}, Legendre and split.
\end{enumerate}

%=============================================================================
\section{The model}
\label{sec:model}

\paragraph{The split of the counterterm.}  The label $u$ of a particle is the argument of the target energy (\texttt{pa.s}, as in \texttt{SPH\_RENORM=1}).  With
the code's \texttt{ConstantCs2EOS}, $p=A\,T_{\rm rel}^{\,n}$, $n=1+1/c_s^2$, and with $(T,v)$ the target state of $(u,\bm m,V)$, the counterterm of
\nolinkurl{legendre_plan.tex} is
\st
\delta\ln Z_1=\Delta s^{\rm lab}(u)+\Delta s^{\rm dyn}(T,v,V),\qquad \mathcal N=1-2c_s^2 ,
\stp
\st
\Delta s^{\rm lab}=\big(\tfrac52-\mathcal N\big)\ln u+{\rm const},\qquad
\Delta s^{\rm dyn}=\mathcal N\ln\frac{\gamma V}{\bar V}+\tfrac12\ln\big(1-c_s^2v^2\big) .
\label{split}
\stp
The first is a function of the label alone, so a factor of it in the measure is conserved by the ideal step (the labels are its Casimirs).  The second depends on
$\bm m$ and $V$, so it goes into $H$.  $\bar V$ is a constant reference volume; its logarithm moves a constant between the two (\Sect{sec:later}).

\paragraph{Measure and Hamiltonian.}  The model is
\begin{keyeq}
\st
\begin{gathered}
d\mu=\prod_a e^{\,u_a-\Delta s^{\rm lab}(u_a)}J^{\rm lab}(u_a)\,du_a\,d\bm x_a\,d\bm m_a,\qquad J^{\rm lab}=1-\frac{\tfrac52-\mathcal N}{u},\\
H_a=e_{\rm t}(u_a,\bm m_a,V_a)+E^{\rm dyn}_a,\\
E^{\rm dyn}_a=T_a\Big(\Delta s^{\rm dyn}_a+\frac{\partial\Delta s^{\rm dyn}_a}{\partial u}\Big|_{\bm m,V}\Big) .
\end{gathered}
\label{model}
\stp
\end{keyeq}
The derivative in $E^{\rm dyn}$ is $-\ln J^{\rm dyn}$ to lowest order: moving the weight $e^{-\Delta s^{\rm dyn}}$ into $H$ also moves the Jacobian of its change of
variables.  The measure vanishes below the label boundary
\st
u_b=\tfrac52-\mathcal N\qquad(2.5\ \text{at}\ c_s^2=\tfrac12,\ 2\ \text{at}\ c_s^2=\tfrac14) .
\stp
At $c_s^2=1/2$ the measure is that of \texttt{SPH\_RENORM=1} ($\alpha/2=5/2$) and $E^{\rm dyn}$ is the sound term alone.

\paragraph{The target state.}  In $(T,v)$ the target is explicit:
\st
u=nAV\gamma^nT^{n-1},\qquad |\bm m|=nAV\gamma^{n+2}T^nv,\qquad e_{\rm t}=AV(\gamma T)^n\big(n\gamma^2-1\big),\qquad P_{\rm t}=A(\gamma T)^n ,
\stp
and the Jacobian entries, with $u_T\equiv\partial u/\partial T$ and so on, and $\det=u_Tm_v-u_v^2$ (the Maxwell relation $m_T=u_v$):
\st
\begin{gathered}
u_T=n(n-1)AV\gamma^nT^{n-2},\qquad u_v=m_T=n^2AV\gamma^{n+2}T^{n-1}v,\\
m_v=nAV\gamma^{n+2}T^n\big(1+(n+2)\gamma^2v^2\big) .
\end{gathered}
\stp
A function $g(T,v,V)$ is differentiated at fixed particle variables by
\st
\frac{\partial g}{\partial u}\Big|_{\bm m,V}=\frac{g_Tm_v-g_vu_v}{\det},\qquad
\frac{\partial g}{\partial|\bm m|}\Big|_{u,V}=\frac{g_vu_T-g_Tu_v}{\det},\qquad
\frac{\partial g}{\partial V}\Big|_{u,\bm m}=g_V-\frac{\partial g}{\partial u}\,\frac uV-\frac{\partial g}{\partial|\bm m|}\,\frac{|\bm m|}V ,
\label{chain}
\stp
the last since $u$ and $|\bm m|$ are proportional to $V$ at fixed $(T,v)$.

\paragraph{What the forward map fills.}  The fields of the $n$-scheme (as in \nolinkurl{legendre_plan.tex}, ``What the code needs from a solve''):
\st
\begin{gathered}
\texttt{pa.T}\leftarrow\frac{\partial H}{\partial u}=T+\frac{\partial E^{\rm dyn}}{\partial u},\qquad
\dot{\bm x}=\frac{\partial H}{\partial\bm m}=\bm v+\frac{\partial E^{\rm dyn}}{\partial|\bm m|}\,\hat{\bm m},\\
\texttt{pa.p}\leftarrow P_{\rm t}-\frac{\partial E^{\rm dyn}}{\partial V},\qquad H=e_{\rm t}+E^{\rm dyn} .
\end{gathered}
\label{fields}
\stp
$E^{\rm dyn}$ already contains one derivative of $\Delta s^{\rm dyn}$, so these need second derivatives of $\Delta s^{\rm dyn}$ and of the target state: closed form in
$(T,v)$, written once by hand (there is no automatic differentiation in the project) and checked by finite differences.

%=============================================================================
\section{The single-particle check}
\label{sec:check}

\nolinkurl{examples/renorm_quadrature_split.py}, the exact quadrature of $Z_1(T,\bm v,V)$ with the label gas (\nolinkurl{partition.tex}, ``Step 5 for any $c_s^2$'',
for the full table).  Relative deviation of the energy from the target per particle:
\begin{center}
\small
\begin{tabular}{ccc cc}
\toprule
$c_s^2$ & $s_a$ & $v$ & Step 5 & Legendre \\
\midrule
$1/2$ & $16$ & $0$ & $+4.1\times10^{-2}$ & $+8.5\times10^{-2}$ \\
$1/2$ & $54$ & $0.8$ & $+1.3\times10^{-3}$ & $+7.4\times10^{-3}$ \\
$1/2$ & $541$ & $0.8$ & $+9.7\times10^{-6}$ & $+6.7\times10^{-5}$ \\
$1/4$ & $16$ & $0.8$ & $+1.1\times10^{-1}$ & $+2.1\times10^{-1}$ \\
$1/4$ & $54$ & $0.8$ & $+7.6\times10^{-5}$ & $+1.6\times10^{-2}$ \\
$1/4$ & $541$ & $0.8$ & $-8.6\times10^{-6}$ & $+1.4\times10^{-4}$ \\
\bottomrule
\end{tabular}
\end{center}
The momentum and the pressure behave the same way, and $R$ is of the same size for the two.  Without $E^{\rm dyn}$ the pressure at $c_s^2=1/4$ is off by the predicted
$\mathcal N/\sigma$.  The velocity $|\partial H/\partial\bm m|$ stays below $1$ on the whole grid (largest $0.997$).  At $s_a=16$, $3$--$5\%$ of the weight lies within
$3$ of the label boundary.

%=============================================================================
\section{The work, in order}
\label{sec:work}

Branch \texttt{split} (\Sect{sec:nersc}).  It builds on the restructuring done on \texttt{legendre}: every particle-level map goes through
\texttt{particle\_forward} and \texttt{particle\_inverse} (\nolinkurl{src/pointwise.jl}), the votes through \texttt{vote\_entropy}
(\nolinkurl{src/stochastic.jl}), with the method chosen at load time by \texttt{SPH\_RENORM}.  Step 5 is a third method behind the same functions.

\subsection{A. The decision before the code}

The term $\partial\Delta s^{\rm dyn}/\partial u$ in $E^{\rm dyn}$ is of relative order $1/u$ and costs one more order of derivatives in \Eq{fields}.  Run the quadrature
with and without it (a switch in the script), at $c_s^2=1/2$, $1/4$ and $s_a=16$, $54$, $541$.  Keep it only if it changes the residual at $s_a\ge54$.

\subsection{B. The maps: \texorpdfstring{\nolinkurl{src/renorm_split.jl}}{src/renorm\_split.jl}}

\begin{algobox}{Forward map: $(u,\bm m,V)\to$ the particle's state}
\begin{algorithmic}[1]
\State Solve the target as now (\texttt{target\_forward}, closed form for \texttt{ConstantCs2EOS}): $T$, $\bm v$, $\gamma$, $e_{\rm t}$, $P_{\rm t}$.
\State From $(T,v)$ evaluate $\Delta s^{\rm dyn}$, its derivatives, and the Jacobian entries.
\State Form $E^{\rm dyn}$ and its derivatives in $u$, $|\bm m|$, $V$ by \Eq{chain}.
\State Fill the fields of \Eq{fields}.
\State If $\partial H/\partial u\le0$ or $|\partial H/\partial\bm m|\ge1$, return a state with no Hamiltonian (as \texttt{SPH\_RENORM=legendre} does; counted).
\end{algorithmic}
\end{algobox}

\begin{algobox}{Inverse map for the votes: $(e,\bm m,V)\to$ the label $u$ and the weight}
\begin{algorithmic}[1]
\State Start from the target inverse at the shifted energy: $u_0=U_{\rm t}\big(e-E^{\rm dyn}_a,\bm m,V\big)$, with $E^{\rm dyn}_a$ of the particle's current state.
\State Newton in $u$ alone: $u\leftarrow u-\big(H(u,\bm m,V)-e\big)\big/(\partial H/\partial u)$, each $H$ by the forward map; stop at a relative residual of
  $10^{-14}$.  Each step gains a factor of order $E^{\rm dyn}/(Tu)$, so two or three steps.
\State If $u\le u_b$, or the forward map has no state: the weight is $-\infty$.
\State Otherwise the weight with the Liouville measure:
  \st
  u-\Delta s^{\rm lab}(u)+\ln J^{\rm lab}(u)-\ln\frac{\partial H}{\partial u} .
  \stp
\end{algorithmic}
\end{algobox}

\subsection{C. The switch}

\texttt{SPH\_RENORM=split}, read as a Symbol as \texttt{legendre} is (\nolinkurl{src/sph.jl}), choosing at load time \texttt{particle\_forward},
\texttt{particle\_inverse} and \texttt{vote\_entropy}, so that without it the code is exactly that of before.  \texttt{SPH\_RENORM\_VBAR} sets $\bar V$; the
default is the mean volume of the initial state, recorded in the run record.  Requirements as for \texttt{legendre} (\texttt{check\_legendre}): the energy
Hamiltonian, the $n$-scheme, $d=2$, the Liouville measure, not the grid step.  Hybrid Monte Carlo rejects a trajectory that leaves the domain of $H$, as now.

\subsection{D. Tests: \texorpdfstring{\nolinkurl{test/testsplit.jl}}{test/testsplit.jl}}

\begin{enumerate}
\item The split \Eq{split} adds up to the counterterm of \nolinkurl{legendre_plan.tex} to $10^{-14}$.
\item The identities of \Eq{fields} by central differences of $H$ to $10^{-9}$; the inverse of the forward map is the identity to $10^{-13}$.
\item Reference values ($H$, $\partial H/\partial u$, $\dot{\bm x}$, $P$) for a few states at both $c_s^2$, from the quadrature script.
\item With $\mathcal N=0$ and the sound term off, the votes equal those of \texttt{SPH\_RENORM=1} to $10^{-12}$.
\item Hybrid Monte Carlo: an energy error $\propto\delta t^2$, and acceptance.
\item \texttt{test/run.sh all} with the switch off: every test passes and the regression is bit-identical (the known \texttt{hast\_lq\_*} exception).
\end{enumerate}

\subsection{E. Speed and smoke runs}

Time per update of \nolinkurl{examples/eos_measure.jl} at $h/a=2$ ($N=360$) for the tree level, \texttt{legendre} and \texttt{split}; the forward map should cost little
more than the tree level's.  Smoke runs at rest, $L=0.45$, $h/a=2$, $s_a=54$, two seeds, both $c_s^2$, against the numbers in \nolinkurl{legendre_plan.tex}, ``Smoke
runs''; and one dynamical run (\nolinkurl{examples/stability.jl}, $t=20$) for the drift of $H$.

\subsection{F. The campaign \texttt{eos\_split}}

The lines of \texttt{eos\_legendre} with \texttt{SPH\_RENORM=split}: $c_s^2\in\{1/2,1/4\}$, $s_a\in\{16,54,541\}$, $h/a\in\{2,3\}$, $v\in\{0,0.8\}$, 8 seeds of 300
updates, $192$ tasks; the tree-level and \texttt{SPH\_RENORM=1} lines are reused.  Cost, from the tree level's times in \nolinkurl{figs/legendre_table.csv} (2 threads:
about $0.12$ h at $h/a=2$, $0.45$ h at $h/a=3$) and a factor $1.3$:
\st
96\times(0.12+0.45)\,{\rm h}\times1.3\times\frac{2}{128}\approx1.1\ \text{node-hours},
\stp
shared QOS, after the rerun of \texttt{eos\_legendre} is done.  The analysis is that of \nolinkurl{examples/eos_legendre.py}, with a third Hamiltonian.

%=============================================================================
\section{Results}
\label{sec:results}

\subsection{A. The term \texorpdfstring{$\partial\Delta s^{\rm dyn}/\partial u$}{dDs/du}}

\nolinkurl{examples/renorm_quadrature_split.py} has the model \texttt{step5-du0}, $E^{\rm dyn}=T\Delta s^{\rm dyn}$ without the derivative (chosen with
\texttt{SR\_MODELS}).  Relative deviation of the energy from the target per particle, with the term / without it:
\begin{center}
\small
\begin{tabular}{cc ccc}
\toprule
$c_s^2$ & $v$ & $s_a=54$ & $100$ & $200$ \\
\midrule
$1/2$ & $0$   & $3.1\times10^{-4}$ / $2.4\times10^{-4}$ & $5.5\times10^{-5}$ / $4.6\times10^{-5}$ & $9.8\times10^{-6}$ / $8.6\times10^{-6}$ \\
$1/2$ & $0.8$ & $1.3\times10^{-3}$ / $9.3\times10^{-4}$ & $2.9\times10^{-4}$ / $1.9\times10^{-4}$ & $7.1\times10^{-5}$ / $4.5\times10^{-5}$ \\
$1/4$ & $0$   & $4.7\times10^{-4}$ / $5.4\times10^{-4}$ & $1.0\times10^{-4}$ / $1.1\times10^{-4}$ & $2.4\times10^{-5}$ / $2.5\times10^{-5}$ \\
$1/4$ & $0.8$ & $7.6\times10^{-5}$ / $6.4\times10^{-4}$ & $-2.4\times10^{-4}$ / $-8.5\times10^{-5}$ & $-6.5\times10^{-5}$ / $-2.8\times10^{-5}$ \\
\bottomrule
\end{tabular}
\end{center}
Both fall as $1/s_a^2$; the term changes the coefficient by up to a factor of about two, either way.  Without it the residual is smaller at $c_s^2=1/2$
everywhere and at $c_s^2=1/4$ for $s_a\ge100$; with it, eight times smaller at $c_s^2=1/4$, $s_a=54$, $v=0.8$, where its residual changes sign between
$s_a=54$ and $100$.  The momentum and the pressure behave alike, and at $s_a=541$ both are below $2\times10^{-5}$.  \emph{Decision} (with the owner): the
term is dropped,
\st
E^{\rm dyn}_a=T_a\,\Delta s^{\rm dyn}_a ,
\stp
so the forward map needs only first derivatives of the target state.  The reference curves of the campaign (\nolinkurl{docs/figs/split_quadrature.csv})
are those of this model.

\subsection{B--D. The code and the tests}

STATUS-CODE

%=============================================================================
\section{Later}
\label{sec:later}

\paragraph{The reference volume, and changing it between steps.}  Changing $\bar V$ moves a function of the label alone between the two parts,
\st
\delta\Delta s^{\rm dyn}=-\delta\Delta s^{\rm lab}\equiv\delta(u) ,
\stp
and with the relabeling $u_a\to u_a-\delta(u_a)$ at the change the Hamiltonian and the weight are continuous up to $O(T\delta^2/u)$.  So $\bar V$ may follow the
volume between steps in an expanding fluid and keep $\Delta s^{\rm dyn}$ small; the error of many small changes vanishes as they are refined.  To be checked with the
single particle (draw at one $\bar V$, relabel, compare with the ensemble at the other) before it is used.

\paragraph{A variable speed of sound.}  With $x=u/(\gamma V)$ the rest-frame entropy density,
\st
\delta\ln Z_1=\tfrac32\ln(2\pi u)+f(x)+\ln(\gamma V)+\tfrac12\ln\big(1-c_s^2v^2\big)+1,\qquad f(x)=2\ln T_{\rm rel}(x)-\tfrac12\ln c_s^2\big(T_{\rm rel}(x)\big),
\stp
and the split takes $f(u/\bar V)$ into $\Delta s^{\rm lab}$.  The equation of state must then supply $c_s^2$ and $dc_s^2/dT$, and the label boundary is found by a
solve.  First in the quadrature script, with a crossover equation of state.

\paragraph{Open.}  Whether the $3$--$5\%$ of the weight near the label boundary at $s_a=16$ makes the votes or the heat bath slow; and the correction of the
positions, $\propto1/(\bar nh^2)$, which no single-particle counterterm removes.

\end{document}
